Having characterised structural aliasing theoretically and empirically, the natural next question is whether it can be mitigated, and whether any proposed remedy actually addresses the root cause rather than masking its symptoms.
This section investigates one principled mitigation strategy — \emph{randomised phase-offset patching} — and analyses its effectiveness through the theoretical lens of Sections~\ref{sec:one}–\ref{sec:two} before reporting the corresponding experimental results.

\subsection{Candidate Mitigation: Randomised Phase-Offset Patching}

The phase-locking degeneracy derived in Proposition~1 (Section~\ref{sec:one}) arises because the patch grid starts at a fixed temporal anchor: all tokens for a given input are extracted at positions $\{kS\}_{k \geq 0}$.
A natural perturbation is to draw a uniform random offset $\delta \sim \mathcal{U}\{0, S-1\}$ at inference time and apply the tokeniser starting at position $\delta$ rather than 0.
By shifting the grid, the phase of the signal relative to the patch boundaries changes, and the strict identity condition $x[n] = x[n+S]$ for all $n$ in the patch becomes less likely to hold globally.

\subsection{Theoretical Analysis}

Formally, let $x[n] = A\cos(2\pi f n / f_s + \phi)$ with $f = m f_s / S$.
Under the original grid (offset 0), all patches are identical as shown in Section~\ref{sec:one}.
Under a random offset $\delta$, the $k$-th patch starts at $kS + \delta$, and the patch values become $x[kS + \delta + j] = A\cos(2\pi m(kS + \delta + j)/S + \phi)$.
Since $2\pi m k S / S = 2\pi m k$ is always a multiple of $2\pi$, the term $kS$ vanishes modulo the signal period, and the patches remain identical regardless of $\delta$.
This is a key negative result: \emph{randomised phase-offset patching does not break the degeneracy at exact harmonic frequencies}, because the phase shift introduced by $\delta$ is the same for all patches.
The mitigation may nonetheless reduce the severity of near-harmonic degradation, where the degeneracy is soft rather than exact.

\subsection{Alternative Mitigations and Residual Risks}

\noindent\textbf{[To be completed.]}
This subsection will survey alternative strategies — learnable patch boundaries, multi-resolution tokenisation, and frequency-aware positional encodings — and provide a theoretical assessment of which structural aliasing modes each strategy does and does not address, contextualising the residual risks.
